\section{Daten}\label{sec:daten}

\subsubsection*{Quelle und Anreicherung}

Die Daten stammen aus dem Bestand von \emph{Energy-Charts} des Fraunhofer ISE \cite{EnergyCharts},
der für die Lehrveranstaltung als PostgreSQL-Datenbank mit PostgREST-Schnittstelle
\cite{PostgREST} unter \url{https://dbs.informatik.uni-halle.de/sciencedata} bereitsteht. Genutzt
und miteinander verbunden werden \texttt{energycharts\_totalpower} (Erzeugungsleistung je
Energieträger in 15-Minuten-Auflösung, gefiltert auf 2024 über \texttt{unix\_seconds}) und
\texttt{energycharts\_price} (Börsenpreise, gefiltert auf \texttt{market\_id = "DE-LU"}).

Die Verbindung beider Quellen ist bereits eine inhaltliche Anreicherung: Erst die Kopplung von
physikalischer Erzeugung und ökonomischem Preis über denselben Zeitstempel macht F2 und F4
beantwortbar. Eine weitere Anreicherung mit Wetterdaten des DWD (Globalstrahlung,
Windgeschwindigkeit) haben wir geprüft und bewusst nicht umgesetzt: Sie würde die Erklärung von der
Wirkung (eingespeiste Leistung) auf die Ursache (Wetter) verschieben und damit eine zweite,
meteorologische Aufgabenklasse eröffnen, für die die drei gewählten Ansichten nicht ausgelegt sind.
Abschnitt~\ref{sec:ausblick} greift den Punkt als Erweiterung auf.

\subsubsection*{Technische Bereitstellung}

Der Zugriff auf die Quelle erfolgt einmalig über das Skript
\texttt{scripts/fetch\_energycharts.py}: Basic Authentication am Endpunkt \texttt{/token}, danach
seitenweiser Abruf über \texttt{/rpc/user\_table\_get} ($500$ Zeilen je Anfrage). Zwei
Besonderheiten der Schnittstelle waren wesentlich und sind im Skript berücksichtigt: Numerische
Bereichsbedingungen (Zeitstempel) funktionieren zuverlässig über das Feld \texttt{where\_} mit den
Operatoren \texttt{>=} und \texttt{<}; Textgleichheit (\texttt{country\_id}, \texttt{market\_id})
dagegen über das Feld \texttt{filters} als Objekt, z.\,B.\ \texttt{\{"market\_id": "DE-LU"\}}.

Ergebnis ist die Datei \texttt{data/energycharts\_de\_2024.json} ($3{,}4$\,MB, UTF-8). Sie wird von
der Anwendung zur Laufzeit per HTTP nachgeladen: \texttt{index.html} führt ein \texttt{fetch()} auf
denselben Ursprung aus und übergibt das Ergebnis als Flag an \texttt{Elm.Main.init}. Die Anwendung
setzt also keine lokal gespeicherten Daten voraus und funktioniert unverändert, sobald HTML,
JavaScript und JSON gemeinsam ausgeliefert werden (etwa über GitLab-Pages).

Die Trennung von Abholen und Ausliefern hat drei Gründe: Zugangsdaten dürfen nicht in eine
clientseitige Anwendung gelangen; die Abgabe wäre sonst von Tokengültigkeit und
Serververfügbarkeit abhängig; und die Rohdaten sind mit über $35\,000$ Viertelstundenzeilen und
mehreren Dutzend Spalten für eine interaktive Darstellung zu breit. Der Preis dafür ist der
Verzicht auf eine Live-Aktualisierung, den Abschnitt~\ref{sec:ausblick} diskutiert.

\subsubsection*{Datenvorverarbeitung}

\emph{1. Normierung der Einheiten.} Die Quellspalten tragen den Suffix \texttt{\_in\_gw}, führen die
deutschen Werte aber in MW-Skala (Solar-Mittagswerte um $43\,000$). Das Skript teilt daher durch
$1000$. Ohne diesen Schritt wären alle Achsen um drei Größenordnungen falsch beschriftet; der
Fehler fiel erst beim Vergleich mit bekannten Größenordnungen des deutschen Systems auf.

\emph{2. Aggregation auf Stunden.} Aus vier Viertelstundenwerten wird der arithmetische Mittelwert
gebildet. Der Börsenpreis ist im betrachteten Jahr ein Stundenprodukt, die Erzeugung liegt in
Viertelstunden vor. Die Stunde ist damit die feinste Auflösung, in der beide Größen gemeinsam
definiert sind. Eine Kopplung auf Viertelstundenebene würde den Preiswert dreimal künstlich
wiederholen und Scheingenauigkeit erzeugen. Der Mittelwert ist hier der richtige Aggregator, weil
Leistung über die Stunde integriert der Arbeit entspricht.

\emph{3. Rekonstruktion der Last.} Das Feld \texttt{load\_in\_gw} ist in der verwendeten
\texttt{totalpower}-Sicht für Deutschland durchgehend leer. Da die Last die Bezugsgröße jedes
Anteilswertes ist, wird sie rekonstruiert:
\[
  \mathrm{Last} \;=\; \mathrm{Solar} + \mathrm{Wind}_{\text{onshore}} + \mathrm{Wind}_{\text{offshore}}
                      + \mathrm{Residuallast}.
\]
Die Identität folgt aus der Definition der Residuallast als \enquote{Last abzüglich
dargebotsabhängiger Einspeisung}. Sie ist eine Näherung, weil Biomasse, Wasserkraft und Geothermie
nicht abgezogen werden. Diese Träger zählen in der Quelle nicht zur Residuallastdefinition. Die
Konsequenz wird in Schritt~5 sichtbar.

\emph{4. Abgeleitete Attribute.} Je Stunde entstehen die erneuerbare Summe (Solar, Wind on-/offshore,
Biomasse, Laufwasser, Speicherwasser, Geothermie), die fossile Summe (Braunkohle, Steinkohle, Öl,
Kokereigas, Erdgas) und der Erneuerbarenanteil als $100 \cdot \mathrm{EE}/\mathrm{Last}$. Je
Kalendertag folgen daraus Mittelwerte, das Lastmaximum, Solar- und Windanteil an der mittleren Last
sowie die Zahl der Stunden mit negativem Preis. Die Tagesebene ist kein Komfort, sondern die
Analyseeinheit für parallele Koordinaten und Detailpanel: Eine einzelne extreme Stunde sagt wenig,
ein Tagesprofil charakterisiert eine Wetterlage.

\emph{5. Werte über 100\,\%.} In $385$ der $8\,690$ Stunden liegt der Erneuerbarenanteil über
$100\,\%$, im Maximum bei $138{,}6\,\%$. Das ist kein Fehler, sondern Folge von
Überschusserzeugung, die exportiert oder eingespeichert wird, und der Näherung aus Schritt~3. Diese
Werte werden nicht auf $100\,\%$ gekappt. Kappen würde genau die Extremsituationen unsichtbar
machen, nach denen F2 sucht. Stattdessen läuft die Farbskala der Heatmap bis $120\,\%$ und sättigt
darüber (Abschnitt~\ref{sec:vis2}).

Die Vorverarbeitung liegt in Python und nicht in Elm. Das ist eine bewusste Abweichung von der
Empfehlung der Aufgabenstellung. Sie gehört zum Abruf aus einer authentifizierten Datenbank, den
die Anwendung aus Sicherheitsgründen nicht selbst ausführen soll. Alles, was nach dem Laden
geschieht, ist in Elm implementiert: Monatsfilterung, Monatsaggregation, Achsenskalierung,
Brush-Auswertung und Kennzahlen. Diese Schritte reagieren unmittelbar auf geänderte Rohdaten; ein
neuer Jahrgang in der JSON-Datei erfordert keine Codeänderung.

\subsubsection*{Umfang, Lücken und kritische Eignung}

\begin{table}[H]
\centering\small
\caption{Umfang des erzeugten Datensatzes}\label{tab:umfang}
\begin{tabular}{ll}
\toprule
Eigenschaft & Wert \\
\midrule
Land / Gebotszone & Deutschland / DE-LU \\
Zeitraum & 1.1.2024 bis 31.12.2024, Zeitstempel in UTC \\
Stundenwerte & $8\,690$ (von $8\,784$ möglichen) \\
Tagesprofile & $364$ (von $366$ möglichen) \\
Stunden mit Preisangabe & $8\,690$ (vollständig) \\
Stunden mit negativem Preis & $441$ ($5{,}1\,\%$) \\
Stunden mit EE-Anteil über $100\,\%$ & $385$ ($4{,}4\,\%$) \\
Mittlerer EE-Anteil / mittlerer Preis & $55{,}1\,\%$ / $78{,}6$\,EUR/MWh \\
Spannweite EE-Anteil (Stunde) & $12{,}0\,\%$ bis $138{,}6\,\%$ \\
Spannweite Preis (Stunde) & $-135{,}45$ bis $936{,}28$\,EUR/MWh \\
\bottomrule
\end{tabular}
\end{table}

Dem Datensatz fehlen $94$ Stunden, die in der Quelle nicht enthalten sind: zwei zusammenhängende
Ausfälle von je $47$ Stunden (23.05.\ 13:00 bis 25.05.\ 11:59 und 12.10.\ 13:00 bis 14.10.\ 11:59,
jeweils UTC). Der 24.05.\ und der 13.10.\ fehlen dadurch vollständig, vier weitere Tage sind nur mit
12 oder 13 Stunden besetzt und ihre Tagesmittel entsprechend schwächer gestützt. Die Lücken werden
nicht interpoliert: In einer pixelorientierten Darstellung wären interpolierte Werte visuell nicht
von Messwerten zu unterscheiden. Damit ginge genau die Eigenschaft verloren, die eine dichte
Darstellung wertvoll macht, nämlich dass jede Zelle für eine Beobachtung steht. Stattdessen zeichnet
die Heatmap zuerst ein vollständiges Raster neutraler grauer Zellen und legt die Datenzellen
darüber; wo Daten fehlen, bleibt Grau sichtbar und ist in der Legende als \enquote{keine Daten}
beschriftet.

Für die Fragestellungen sind die Daten gut geeignet: Die Preisreihe ist lückenlos, die zeitliche
Auflösung liegt deutlich unter der kürzesten interessierenden Struktur (dem Tagesgang), und
Erzeugung und Preis stammen aus demselben Marktgebiet. Drei Einschränkungen sind zu nennen. Erstens
ist die Last rekonstruiert; alle Anteilswerte sind Näherungen, deren systematische Richtung bekannt
ist. Sie fallen tendenziell zu hoch aus. Zweitens ist die Vorzeichenkonvention von
\texttt{cross\_border\_electricity\_trading\_in\_gw} in der Quelle nicht dokumentiert; der Wert wird
deshalb neutral als Nettohandelswert geführt. Die Daten stützen die Lesart \enquote{negativ =
Export}: Am Tag mit dem höchsten Erneuerbarenanteil (4.2., $95{,}9\,\%$) beträgt er $-13{,}1$\,GW, in
der Dunkelflautenstunde am 6.11.\ dagegen $+13{,}4$\,GW. Drittens zeigen die Daten Koinzidenz, nicht
Kausalität; die Anwendung ist ein Werkzeug zur Hypothesenbildung im Sinne von \cite{Keim2008}, kein
Erklärungsmodell.
